\documentclass{article}
% Source: Topic_07_Teacher_Edition.pdf, PDF pages 81-82 (printed pages 391A-391B): Topic Assessment Form A with the TE's marked answers, Form B with marked answers, Item Analysis for both forms. No extra text source for this topic; transcribed from the 250 dpi page renders and 600 dpi crops of each student page.
\usepackage{amsmath}
\usepackage{amssymb}
\newenvironment{assessment-meta}{}{}   % topic, title, forms, item count, printed pages, sources
\newenvironment{item-analysis}{}{}     % Savvas's Item / DOK / Skills Review and Practice table (same for both forms)
\newenvironment{ta-item}[2]{}{}        % #1 = item number, #2 = form (A or B)
\newcommand{\answer}[1]{}              % the TE's marked answer, inside a ta-item
\newcommand{\placeholder}[2]{}         % #1 = type (graph, table, figure, number-line), #2 = complete description
\newcommand{\uncertain}[2]{#1}         % #1 = best reading, #2 = alternatives and why
\newcommand{\te}[1]{}                  % transcriber's note (a discrepancy, a publisher misprint, a computed check)
\begin{document}
\begin{assessment-meta}
topic: 7
title: Trigonometric Functions
forms: A and B (parallel item for item)
items: 18 per form
printed: 391A-391B
sources: Topic_07_Teacher_Edition.pdf, PDF pages 81-82 (printed 391A-391B); no extra text source
\end{assessment-meta}
\te{Page 391A (PDF 81) running head: "TOPIC 7 Trigonometric Functions", banner "TOPIC ASSESSMENT Form A", "AVAILABLE ONLINE"; footer "TOPIC 7 | 391A | TOPIC ASSESSMENT". Page 391B (PDF 82): same running head, banner "TOPIC ASSESSMENT Form B"; footer "TOPIC 7 | 391B | TOPIC ASSESSMENT". Student pages headed "7 Topic Assessment Form A" / "7 Topic Assessment Form B", with Name line and the enVision Algebra 2 / SavvasRealize.com badge; student footers "enVision(R) Algebra 2 | 1 of 2 | Assessment Sourcebook" and "... 2 of 2 ...". Running heads, student headers and footers agree; no topic misprint on these pages.}
\te{TE text under Form A: "Topic Assessment Assess students' understanding of topic concepts and skills using the Topic Assessment found at SavvasRealize.com. These auto-scored online assessments provide students with a breadth of technology-enhanced item types. There are two versions of the Topic Assessment, Form A and Form B. These two versions, available in print and at SavvasRealize.com, are parallel tests that assess the same content item for item. The Item Analysis chart on the next page can be used for both versions."}
% Source page 82: 391B Item Analysis
\begin{item-analysis}
Item & DOK & Skills Review and Practice \\
1 & 2 & F86 \\
2 & 2 & F90 \\
3 & 2 & F86 \\
4 & 2 & F84 \\
5 & 2 & F89 \\
6 & 2 & F88, F90 \\
7 & 2 & F88, F90 \\
8 & 2 & F88 \\
9 & 2 & F84 \\
10 & 1 & F90 \\
11 & 2 & F31 \\
12 & 1 & F87 \\
13 & 2 & F84 \\
14 & 2 & F85 \\
15 & 2 & F89 \\
16 & 2 & F86 \\
17 & 2 & F87, F90 \\
18 & 2 & F31, F87 \\
\end{item-analysis}
% Source page 81: 391A Topic Assessment Form A (student page 1 of 2)
\begin{ta-item}{1}{A}
Given (\sin\theta=\frac{12}{13}), what is (\sec\theta)?
\frac{[ ]}{[ ]}
\answer{1A. numerator box 13, denominator box 5, i.e. (\frac{13}{5})}
\te{The item names no quadrant or sign restriction; the printed key takes the positive value (an acute angle, Quadrant I). Codex review 2026-10-31.}
\end{ta-item}
\begin{ta-item}{2}{A}
The graph of (f) is given. How are (f) and (g(x)=2\cos(\pi x)) different?
\placeholder{graph}{Part of the prompt. Square grid, gridlines every 1 unit; x from about -1 to 8, y from about -3 to 3; axes with arrowheads labeled x and y, origin O; x labels 2, 4, 6; y labels 2 and -2. Smooth cosine-shaped curve of amplitude 2: maximum (0,2) on the y-axis, minimum (2,-2), maximum (4,2) marked with a dot and labeled "(4, 2)", minimum (6,-2); crosses the x-axis at 1, 3, 5, 7; arrows at both ends (left end falling from the maximum (0,2) down-left to an arrowhead at about (-0.9, 0.4), between the x-axis and y=1; right end rising past x=7 near y=2). Period 4.}
A. The period of (f) is half as long as the period of (g).
B. The period of (f) is twice as long as the period of (g).
C. The amplitude of (f) is half as great as the amplitude of (g).
D. The amplitude of (f) is twice as great as the amplitude of (g).
\answer{2A. B}
\end{ta-item}
\begin{ta-item}{3}{A}
If (\cot\theta=3), what is (\sin\theta)? Round to the nearest hundredth. [ ]
\answer{3A. 0.32}
\te{The item names no quadrant or sign restriction; the printed key takes the positive value (an acute angle, Quadrant I). Codex review 2026-10-31.}
\end{ta-item}
\begin{ta-item}{4}{A}
What are the sine and cosine of (\frac{4\pi}{3})?
\placeholder{figure}{Six answer tiles in two rows: (-\frac{\sqrt{3}}{2}), (\sqrt{3}), (\frac{\sqrt{3}}{2}); (-\frac{1}{2}), 2, (\frac{1}{2}).}
(\sin\frac{4\pi}{3}=) [ ]
(\cos\frac{4\pi}{3}=) [ ]
\answer{4A. (\sin\frac{4\pi}{3}=-\frac{\sqrt{3}}{2}); (\cos\frac{4\pi}{3}=-\frac{1}{2})}
\end{ta-item}
\begin{ta-item}{5}{A}
Tyler is observing the ascent of a steel beam being lifted by a crane from a position about 250 feet away. What function models the height (h) in feet of the beam as a function of the angle of inclination (\theta) from Tyler's position to the steel beam?
(h=) [ ] [ ] (\theta)
\answer{5A. (h=250\tan\theta): boxes 250, tan}
\end{ta-item}
\begin{ta-item}{6}{A}
Today, high tide measured 7.5 ft and low tide measured 3 ft. What is the amplitude of a cosine function modeling the depth of the water in feet as a function of time in hours? [ ] ft
\answer{6A. 2.25}
\end{ta-item}
\begin{ta-item}{7}{A}
Write a trigonometric function that models the average high temperatures shown in the graph.
\placeholder{graph}{Part of the prompt. Title "Average High Temperature (\textdegree F)"; horizontal axis x, labeled "Hours", labels 0, 6, 12, 18, 24, 30, 36, 42, 48 (gridlines every 3 hours); vertical axis y, window from 0 to 17.5, labels 0, 5, 10, 15 (gridlines every 2.5). Scatter of black dots about every 1.5 hours from x=1.5 to x=48 tracing a cosine wave: near 15 at the start (about 14.5 at x=1.5), minimum about 2.5 at x=12, maximum 15 at x=24, minimum about 2.5 at x=36, back to 15 at x=48. Period 24, midline 8.75, amplitude 6.25.}
\placeholder{figure}{Six answer tiles in two rows: -6.25, 6.25, -8.75; (\cos(\frac{\pi}{12}x)), 8.75, (\sin(\frac{\pi}{12}x)).}
(y=) [ ] [ ] (+) [ ]
\answer{7A. (y=6.25\cos(\frac{\pi}{12}x)+8.75): boxes 6.25, (\cos(\frac{\pi}{12}x)), 8.75}
\end{ta-item}
\begin{ta-item}{8}{A}
Function (f) is a cosine function with period (8\pi), amplitude 2, and a local maximum at (f(0)=-1). What is the equation of the midline of the graph of (f)? (y=) [ ]
\answer{8A. -3}
\end{ta-item}
% Source page 81: 391A Topic Assessment Form A (student page 2 of 2)
\begin{ta-item}{9}{A}
Select all the possible measures for the angle shown.
\placeholder{graph}{Part of the prompt. Coordinate axes with arrowheads, labeled x and y, origin O; no grid or scale. The initial side lies along the positive x-axis (drawn bold with an arrowhead); the terminal ray, with an arrowhead, points down and to the left into Quadrant III, making an angle marked 30\textdegree{} with the negative x-axis (arc between the negative x-axis and the ray).}
A. (-\frac{17\pi}{6})
B. (-210^\circ)
C. (-\frac{2\pi}{3})
D. (\frac{5\pi}{6})
E. (210^\circ)
\answer{9A. A, E}
\end{ta-item}
\begin{ta-item}{10}{A}
What is the phase shift of the graph of (y=-\sin(x+\frac{\pi}{4})+2)?
A. (\frac{\pi}{4}) units left
B. 2 units left
C. (\frac{\pi}{4}) units right
D. 2 units right
\answer{10A. A}
\end{ta-item}
\begin{ta-item}{11}{A}
What is the average rate of change for the function in Item 12, over the interval ([0,2\pi])?
A. -1
B. 0
C. 2
D. (2\pi)
\answer{11A. B}
\te{Item 12 on this form is "If sin(pi/10) = a ..." and has no function to evaluate; the cross-reference is a publisher misprint (likely online item numbering). The printed key 0 agrees with Item 10's function: (y=-\sin(x+\frac{\pi}{4})+2) has (y(2\pi)=y(0)), so the rate is 0. (Item 18's (y=2\sin(\frac{1}{2}x)) also gives 0 on ([0,2\pi]).) Choices are printed in two columns: A and B on the left, C and D on the right.}
\end{ta-item}
\begin{ta-item}{12}{A}
If (\sin\frac{\pi}{10}=a), select all of the expressions that are equivalent to (a).
A. (\sin(\frac{\pi}{10}+\pi))
B. (\sin(\frac{\pi}{10}+2\pi))
C. (\sin(\frac{\pi}{10}+3\pi))
D. (\sin(\frac{\pi}{10}+4\pi))
E. (\sin(\frac{\pi}{10}+5\pi))
\answer{12A. B, D}
\end{ta-item}
\begin{ta-item}{13}{A}
A car is on a Ferris wheel with a radius of 20 feet. To the nearest foot, how far does the car travel over an angle of (\frac{\pi}{3}) radians?
[ ] ft
\answer{13A. 21}
\end{ta-item}
\begin{ta-item}{14}{A}
An angle intercepts an arc of length (\frac{31\pi}{36}) radians on the unit circle. What are a possible positive measure and negative measure for the angle?
A. (335^\circ, -25^\circ)
B. (25^\circ, -335^\circ)
C. (205^\circ, -155^\circ)
D. (155^\circ, -205^\circ)
\answer{14A. D}
\end{ta-item}
\begin{ta-item}{15}{A}
What equation is represented by the graph?
\placeholder{graph}{Part of the prompt. Grid, x gridlines every (\frac{\pi}{4}), y gridlines every 0.5; x from (-\frac{\pi}{2}) to (\frac{5\pi}{2}), y from -2 to 2; axes labeled x and y, origin O; x labels (\pi) and (2\pi), y labels 1 and -1. Dashed vertical asymptotes, with arrowheads at both ends, at (x=-\frac{\pi}{2}), (\frac{\pi}{2}), (\frac{3\pi}{2}), (\frac{5\pi}{2}). Three increasing tangent branches through (0,0), ((\pi,0)), ((2\pi,0)), each passing through height 0.5 one gridline ((\frac{\pi}{4})) right of its zero and leaving the window with arrows at top and bottom next to the asymptotes; the branches flatten near the zeros (vertical stretch (\frac{1}{2})).}
A. (y=\frac{1}{2}\tan x)
B. (y=\frac{1}{2}\cot x)
C. (y=\tan\frac{1}{2}x)
D. (y=\cot\frac{1}{2}x)
\answer{15A. A}
\end{ta-item}
\begin{ta-item}{16}{A}
What is (\cos\theta) if (\sin\theta=-\frac{8}{10}) and (\theta) is in Quadrant III?
A. (\cos\theta=-\frac{6}{10})
B. (\cos\theta=\frac{6}{10})
C. (\cos\theta=-\frac{6}{8})
D. (\cos\theta=\frac{6}{8})
\answer{16A. A}
\end{ta-item}
\begin{ta-item}{17}{A}
What is the frequency of the graph of (y=2\sin(\frac{1}{2}x+2\pi)-3)?
\placeholder{figure}{Six answer tiles in two rows: (\frac{1}{2\pi}), (\frac{1}{4\pi}), (\pi); (2\pi), (4\pi), (\frac{1}{\pi}). One answer box below.}
[ ]
\answer{17A. (\frac{1}{4\pi})}
\end{ta-item}
\begin{ta-item}{18}{A}
What is the average rate of change of the function (y=2\sin(\frac{1}{2}x)) on the interval ([0,\pi])?
A. 2
B. (\frac{\pi}{2})
C. 1
D. (\frac{2}{\pi})
\answer{18A. D}
\end{ta-item}
% Source page 82: 391B Topic Assessment Form B (student page 1 of 2)
\begin{ta-item}{1}{B}
Given (\sin\theta=\frac{24}{25}), what is (\cos\theta)?
\frac{[ ]}{[ ]}
\answer{1B. numerator box 7, denominator box 25, i.e. (\frac{7}{25})}
\te{The item names no quadrant or sign restriction; the printed key takes the positive value (an acute angle, Quadrant I). Codex review 2026-10-31.}
\end{ta-item}
\begin{ta-item}{2}{B}
The graph of (f) is given. How are (f) and (g(x)=2\sin(\pi x)) different?
\placeholder{graph}{Part of the prompt. Square grid, gridlines every 0.5 on both axes; x from -1 to about 4.5, y from -2.5 to 2.5; axes with arrowheads labeled x and y, origin O; x labels 1, 2, 3; y labels 2 and 1. Sine-shaped curve of amplitude 2 and period 1: through (0,0) rising, maxima at x=-0.75, 0.25, 1.25, 2.25, 3.25, 4.25 (y=2), minima at x=-0.25, 0.75, 1.75, 2.75, 3.75 (y=-2); a dot at (1,0) labeled "(1, 0)" by a leader line from below the axis; arrows at both ends (left end falling at about x=-1, right end falling at about x=4.5). A small stray arrowhead sits on the rising part of the curve just below the axis near x=3.9.}
A. The period of (f) is half as long as the period of (g).
B. The period of (f) is twice as long as the period of (g).
C. The amplitude of (f) is half as great as the amplitude of (g).
D. The amplitude of (f) is twice as great as the amplitude of (g).
\answer{2B. A}
\end{ta-item}
\begin{ta-item}{3}{B}
If (\tan\theta=2), what is (\cos\theta)? Round to the nearest hundredth. [ ]
\answer{3B. 0.45}
\te{The item names no quadrant or sign restriction; the printed key takes the positive value (an acute angle, Quadrant I). Codex review 2026-10-31.}
\end{ta-item}
\begin{ta-item}{4}{B}
What are the sine and cosine of (\frac{5\pi}{3})?
\placeholder{figure}{Six answer tiles in two rows: (-\frac{\sqrt{3}}{2}), (\sqrt{3}), (\frac{\sqrt{3}}{2}); (-\frac{1}{2}), 2, (\frac{1}{2}).}
(\sin\frac{5\pi}{3}=) [ ]
(\cos\frac{5\pi}{3}=) [ ]
\answer{4B. (\sin\frac{5\pi}{3}=-\frac{1}{2}); (\cos\frac{5\pi}{3}=-\frac{\sqrt{3}}{2})}
\te{computed: (\sin\frac{5\pi}{3}=-\frac{\sqrt{3}}{2}) and (\cos\frac{5\pi}{3}=\frac{1}{2}) (reference angle (\frac{\pi}{3}), Quadrant IV). The printed key is wrong on both boxes; its values are those of (\frac{7\pi}{6}). Confirmed on the 600 dpi crop.}
\end{ta-item}
\begin{ta-item}{5}{B}
Nicholas observes a box being hoisted onto a ship. Nicholas stands on a dock about 150 feet away. Which function models the height, (h), in feet, of the cargo as a function of the angle of inclination, (\theta), from Nicholas to the box?
(h=) [ ] [ ] (\theta)
\answer{5B. (h=150\tan\theta): boxes 150, tan}
\end{ta-item}
\begin{ta-item}{6}{B}
Today, high tide is measured at 6 ft and low tide is measured at 2.5 ft. What is the amplitude of a cosine function modeling the depth of the water in feet as a function of time in hours? [ ] ft
\answer{6B. 1.75}
\end{ta-item}
\begin{ta-item}{7}{B}
Write a trigonometric function that models the average rainfall shown in the graph?
\placeholder{graph}{Part of the prompt. Title "Average Rainfall (in.)"; horizontal axis x, labeled "Months", labels 0, 3, 6, 9, 12, 15, 18, 21, 24 (gridlines every 1.5); vertical axis y, window from 0 to 7, labels 0, 2, 4, 6 (gridlines every 1); x window from 0 to 24. Black dots at each whole month x=0 to 24 tracing a cosine wave: maximum 6 at x=0, minimum 1 at x=6, maximum 6 at x=12, minimum 1 at x=18, maximum 6 at x=24 (about 3.5 at x=3, 9, 15, 21). Period 12, midline 3.5, amplitude 2.5.}
\placeholder{figure}{Six answer tiles in two rows: -2.5, 2.5, -3.5; (\cos(\frac{\pi}{6}x)), 3.5, (\sin(\frac{\pi}{6}x)).}
(y=) [ ] [ ] (+) [ ]
\answer{7B. (y=2.5\cos(\frac{\pi}{6}x)+3.5): boxes 2.5, (\cos(\frac{\pi}{6}x)), 3.5}
\te{Form B's prompt ends with "?" after "in the graph" (Form A ends with a period); kept as printed.}
\end{ta-item}
\begin{ta-item}{8}{B}
Function (f) is a cosine function with period (8\pi), amplitude 2.5, and a local maximum at (f(0)=6). What is the equation of the midline of the graph of (f)? (y=) [ ]
\answer{8B. 3.5}
\end{ta-item}
% Source page 82: 391B Topic Assessment Form B (student page 2 of 2)
\begin{ta-item}{9}{B}
Select all the possible measures for the angle shown.
\placeholder{graph}{Part of the prompt. Coordinate axes with arrowheads, labeled x and y, origin O; no grid or scale. The initial side lies along the positive x-axis (drawn bold with an arrowhead); the terminal ray, with an arrowhead, points down and to the right into Quadrant IV, making an angle marked 30\textdegree{} with the positive x-axis (arc between the positive x-axis and the ray).}
A. (-\frac{23\pi}{6})
B. (-\frac{13\pi}{6})
C. (-\frac{7\pi}{6})
D. (-330^\circ)
E. (330^\circ)
\answer{9B. B, E}
\end{ta-item}
\begin{ta-item}{10}{B}
What is the phase shift of the graph of (y=-\sin(x-\pi)-3)?
A. (\pi) units left
B. 3 units left
C. (\pi) units right
D. 3 units right
\answer{10B. C}
\end{ta-item}
\begin{ta-item}{11}{B}
What is the average rate of change for the function in Item 12, over the interval ([0,\pi])?
A. (-\pi)
B. -1
C. 0
D. (\pi)
\answer{11B. C}
\te{As on Form A, Item 12 has no function; publisher cross-reference misprint. The key 0 agrees with Item 10's (y=-\sin(x-\pi)-3) ((y(\pi)=y(0)=-3)). Choices printed in two columns: A and B on the left, C and D on the right.}
\end{ta-item}
\begin{ta-item}{12}{B}
If (\sin\frac{\pi}{8}=a), select all of the expressions that are equivalent to (a).
A. (\sin(\frac{\pi}{8}+\pi))
B. (\sin(\frac{\pi}{8}+2\pi))
C. (\sin(\frac{\pi}{8}+4\pi))
D. (\sin(\frac{\pi}{8}+6\pi))
E. (\sin(\frac{\pi}{8}+8\pi))
\answer{12B. B, C, D, E}
\end{ta-item}
\begin{ta-item}{13}{B}
Reagan rides on a playground roundabout with a radius of 2.5 feet. To the nearest foot, how far does Reagan travel over an angle of (\frac{4\pi}{3}) radians? [ ] ft
\answer{13B. 10}
\end{ta-item}
\begin{ta-item}{14}{B}
An intercepts an arc of length (\frac{29\pi}{36}) radians on the unit circle. What are a possible positive measure and negative measure for the angle?
A. (325^\circ, -35^\circ)
B. (35^\circ, -325^\circ)
C. (215^\circ, -145^\circ)
D. (145^\circ, -215^\circ)
\answer{14B. D}
\te{"An intercepts" is printed so: the word "angle" is missing (Form A reads "An angle intercepts"). Kept verbatim.}
\end{ta-item}
\begin{ta-item}{15}{B}
What equation is represented by the graph?
\placeholder{graph}{Part of the prompt. Grid, x gridlines every (\frac{\pi}{4}), y gridlines every 0.5; x from (-\frac{\pi}{2}) to (\frac{5\pi}{2}), y from -2 to 2; axes labeled x and y, origin O; x labels (\pi) and (2\pi), y labels 1 and -1. Dashed vertical asymptotes, with arrowheads at both ends, at (x=-\frac{\pi}{2}), (\frac{\pi}{2}), (\frac{3\pi}{2}), (\frac{5\pi}{2}). Three steep, nearly straight increasing tangent branches through (0,0), ((\pi,0)), ((2\pi,0)), each leaving the window with arrows at (y=2) about (\frac{\pi}{4}) right of its zero and at (y=-2) about (\frac{\pi}{4}) left of it (vertical stretch 2).}
A. (y=2\cot x)
B. (y=2\tan x)
C. (y=\tan 2x)
D. (y=\cot 2x)
\answer{15B. B}
\end{ta-item}
\begin{ta-item}{16}{B}
What is (\cos\theta) if (\sin\theta=-\frac{8}{10}) and (\theta) is in Quadrant IV?
A. (\cos\theta=-\frac{6}{10})
B. (\cos\theta=\frac{6}{10})
C. (\cos\theta=-\frac{6}{8})
D. (\cos\theta=\frac{6}{8})
\answer{16B. B}
\end{ta-item}
\begin{ta-item}{17}{B}
What is the frequency of the graph of (y=4\sin(x+\pi)-2)?
\placeholder{figure}{Six answer tiles in two rows: (\frac{1}{2\pi}), (\frac{1}{4\pi}), (\pi); (2\pi), (4\pi), (\frac{1}{\pi}). One answer box below.}
[ ]
\answer{17B. (\frac{1}{2\pi})}
\end{ta-item}
\begin{ta-item}{18}{B}
What is the average rate of change of the function (y=4\sin(x)) on the interval ([0,\frac{3\pi}{2}])?
A. (-\frac{16}{\pi})
B. (-\frac{8}{3\pi})
C. (-\frac{4}{\pi})
D. (-\frac{2}{\pi})
\answer{18B. B}
\end{ta-item}
\end{document}
